% ECONOMICS_PROBLEM: {"id": "O34-204", "title": "Training-data compensation", "jel_code": "O34", "source_id": "OP-0415"}
% !TeX program = xelatex
\documentclass[11pt,a4paper]{article}
\usepackage[margin=23mm]{geometry}
\usepackage{fontspec}
\setmainfont{Latin Modern Roman}
\usepackage{amsmath,amssymb,mathtools}
\usepackage{xcolor,enumitem,booktabs,longtable,array,xurl,hyperref}
\definecolor{muted}{HTML}{53636C}
\hypersetup{colorlinks=true,urlcolor=blue,linkcolor=blue}
\setlength{\parindent}{0pt}
\setlength{\parskip}{5pt}
\newcommand{\R}{\mathbb R}
\newcommand{\E}{\mathbb E}
\newcommand{\N}{\mathbb N}
\newcommand{\ind}{\mathbf 1}
\newcommand{\Var}{\operatorname{Var}}
\newcommand{\Cov}{\operatorname{Cov}}
\newcommand{\supp}{\operatorname{supp}}
\DeclareMathOperator*{\argmin}{arg\,min}
\DeclareMathOperator*{\argmax}{arg\,max}
\newcommand{\entrymeta}[3]{{\sffamily\small\color{muted}\textbf{\texttt{#1}}\quad|\quad #2\par #3\par}}
\newcommand{\Problemref}[1]{\texttt{#1}}
\newcommand{\JELref}[2]{\texttt{#2} (JEL \texttt{#1})}
\begin{document}
\section*{Training-data compensation}
\textbf{Problem ID:} \texttt{O34-204}\quad \textbf{JEL:} \texttt{O34}\par
% BEGIN ECONOMICS STATEMENT
\label{problem:OP-0415}
{\sffamily\small\color{muted}Original subject group: Innovation and AI economics\par}
\label{definitions:problem:30007697}
\textbf{Audit classification:} Background; proposed extension. The intellectual-property paper is verified background. The proposed compensation mechanism is not a source-authored canonical conjecture; equilibrium selection and maximizer existence were missing.
\subsection*{Definitions and precise research target}
Fix a market containing original-content creators, one model developer and downstream users. Creator \(i\) chooses costly original-content quality \(x_i\geq0\), with real cost \(c_i(x_i)\); the developer obtains licensed training content and produces model quality \(q(x)\), where \(x\) is the vector of creator choices. A compensation mechanism \(r\) specifies admissible training uses, observable quality verification and transfers \(t_i(x)\) to creators. Let \(\mathcal R\) be a stated feasible mechanism class satisfying budget balance, creator participation and the developer's participation constraint. Creator utility is \(t_i(x)-c_i(x_i)\), plus separately specified revenue from nontraining uses. Let \(B(q,x)\) be users' aggregate gross value from the model and original works, and \(K(q)\) the developer's real training and deployment cost. With transfers counted only once, social welfare is
\[
W(r)=B(q(x(r)),x(r))-\sum_i c_i(x_i(r))-K(q(x(r))),
\]
where \(x(r)\) is an equilibrium quality profile under mechanism \(r\). Determine when a welfare-maximizing \(r^\ast\in\arg\max_{r\in\mathcal R}W(r)\) can be implemented using verifiable usage and quality information, and how it changes future content supply. This is a proposed mechanism problem; conclusions depend on specified outside options, production technology and equilibrium selection, and assert no universal legal compensation rule.

\textbf{Definitions, assumptions and mathematical qualifications.} The mechanism class includes licensing acceptance, transfer sources and verifiable quality or usage reports, not just \(t_i(x)\). Budget balance specifies whether payments are funded from user fees, developer revenue or a public subsidy; all included transfers cancel from total welfare. Let creators choose quality in fixed compact intervals and specify the developer pricing/investment game and all outside options. Admit only mechanisms with nonempty equilibrium sets and choose a supported kernel \(\kappa_r\). Replace \(x(r)\) by selected equilibrium outcomes or by integration over \(\kappa_r\). A finite nonempty feasible mechanism menu gives existence of a welfare maximizer; unrestricted \(\mathcal R\) needs regularity or an \(\varepsilon\)-optimal target.
\subsection*{Variants and assumption changes}
\begin{enumerate}
\item \textbf{Editorial, known observable qualities.} With a finite contract menu and known quality technology, compare equilibria and welfare by finite computation; this is not a general open existence theorem.
\item \textbf{Editorial, latent contribution.} Permit unverifiable marginal training contribution and strategic creators. Restrict payments to verifiable usage signals and study implementable welfare ranges and future original-content supply.
\end{enumerate}
\subsection*{References and source audit}
\textbf{Checked source.} NBER Working Paper 32698 (July 2024), indexed abstract; published chapter DOI 10.1086/732853 (2025). The paper develops an economic intellectual-property framework; these exact optimization questions are editorial. Primary metadata and abstract indexed; full official chapter returned HTTP 403. \url{https://www.nber.org/papers/w32698}.
\begin{enumerate}\raggedright
\item\hypertarget{definitions:source:30007697:1}{} Gaétan de Rassenfosse; Adam B. Jaffe; Joel Waldfogel. 2024. \emph{Intellectual Property and Creative Machines}. NBER Working Paper 32698 (July 2024), indexed abstract; published chapter DOI 10.1086/732853 (2025).
\url{https://www.nber.org/papers/w32698}.
DOI: \url{https://doi.org/10.3386/w32698}.
\end{enumerate}

\subsection*{Common conventions: Source mathematical conventions}

These conventions apply to each entry unless explicitly replaced.

\textbf{Models and identification.} A primitive model \(m\) specifies all probability laws, feasible choices, information, equilibrium and measurement rules used by its target. The admissible class \(\mathcal M\) is fixed independently of outcomes. If \(P_O(m)\) is its observable probability law and \(\tau(m)\) the target, its identified set at observable law \(P\) is
\[
 \mathcal I_\tau(P)=\{\tau(m):m\in\mathcal M,\ P_O(m)=P\}.
\]
Point identification means that this set is a singleton. An empty set signals incompatibility of data and assumptions. Coordinate bounds need not describe the joint set. Bounds are sharp only if every retained target is attainable or arbitrarily approachable by compatible models. Finite bounds require explicit restrictions on unobserved quantities; unrestricted confounding, hidden wealth, extrapolation or arbitrary equilibrium selection can yield uninformative sets.

\textbf{Causal interventions.} A potential outcome \(Y_i(p)\) is the outcome under a complete policy regime \(p\), including timing, financing, information and market spillovers. The target population and units are fixed before treatment. With interference, use the complete assignment vector or a justified exposure mapping; individual no-interference notation alone cannot identify an economy-wide intervention. A difference of mean potential outcomes does not require the cross-world joint distribution. Conditioning on post-treatment migration, survival, enrollment or employment changes the estimand and needs separate justification.

\textbf{Statistical statements.} An estimator, confidence set, forecast or test specifies a sampling law, model class and loss/coverage criterion. Uniform \((1-\alpha)\)-coverage means \(\inf_{m\in\mathcal M}\Pr_m\{\tau(m)\in C_n\}\geq1-\alpha\), or an explicitly asymptotic version. The whole parameter space trivially has coverage, so informativeness requires a stated width, risk or power objective. A forecasting improvement is relative to a named benchmark, information set and evaluation population.

\textbf{Equilibrium and optimization.} An equilibrium comprises feasible plans, optimality, beliefs where needed, budget constraints and all market/resource clearing. The primitives or a stated benchmark define each correspondence \(\mathcal E(p;m)\). Nonemptiness is assumed or proved, never inferred just from compact policy sets. A selected-equilibrium target uses a declared selection rule; a robust target quantifies over all permitted equilibria. Use suprema or infima unless attainment is established. An \(\varepsilon\)-optimal policy is feasible and within \(\varepsilon\) of the finite optimum value. Report an empty feasible set separately.

\textbf{Welfare and accounting.} Normative utility, interpersonal normalization, social weights, numeraire, discounting and the use of public receipts are primitives. Behavioral utility need not coincide with the welfare criterion. Household welfare includes ownership income and rents once; adding profits again to compensating variation generally double-counts them. A monetary equivalent is defined by an explicit transfer and price convention, with monotonicity and range conditions. Average effects, mechanism contributions, transfers and welfare losses are different targets. Infinite sums require convergence and dynamic feasible plans require terminal or no-Ponzi conditions.

\textbf{Source versions.} A working-paper date, revision date and journal publication year are distinguished. A DOI for a journal version is not automatically the DOI of an earlier manuscript. Locators refer to the stated version and printed pages where specified. A metadata-only check does not verify a claim about a full-text conclusion. Every entry's source subsection narrows the provenance claim accordingly.
% END ECONOMICS STATEMENT
\end{document}
